\section{Task and Benchmark Construction}

\subsection{Task Definition}

Each instance contains claim text $x_t$, an image $x_v$, a decision cutoff $\tau$, and an evidence collection $E_{\leq\tau}$. The output is $y \in \{\text{supported},\text{refuted},\text{insufficient}\}$. Documents published after $\tau$ are unavailable to the model. The event is the independent unit for data splitting and uncertainty analysis.

\subsection{Candidate Acquisition}

Corpus construction starts with English ClaimReview records about United States events from 2020--2026. GDELT provides discovery metadata, and Common Crawl provides archived page observations. The corpus retains cases with an identifiable claim time, a recoverable image or archive, an unambiguous professional rating, and enough timestamp information to enforce the cutoff. The retrieval pool excludes fact-check domains and documents that contain the final verdict.

\subsection{Label Policy}

Only explicit site ratings are normalized. True or correct maps to supported. False, fake, or miscaptioned maps to refuted. Unproven, unsubstantiated, or no evidence maps to insufficient evidence. Mixed, partially true, satire, opinion, and ambiguous ratings are excluded. The labels come from professional fact-checkers; automatically extracted evidence links are treated as silver evidence.

\subsection{Temporal Partitioning and Release}

Training covers 2020--2023, validation uses 2024, and final testing uses 2025--2026. Events, near-duplicate claims, and images with matching perceptual hashes remain in one partition. The release package contains URLs, cryptographic and perceptual hashes, timestamps, normalized labels, graph annotations, retrieval queries, and split manifests. It does not redistribute copyrighted articles or images.

Figure~\ref{fig:temporal-protocol} shows the order of operations. Group identities are fixed before the year-based allocation. Training data are used to fit model parameters and preprocessing; validation data are reserved for model selection, calibration, and abstention thresholds; the final period is not used for any of those choices.

\begin{figure}[t]
\centering
\resizebox{\columnwidth}{!}{\begin{tikzpicture}[
  font=\sffamily\footnotesize,
  phase/.style={draw, very thick, rounded corners=1.5pt, minimum height=11mm, align=center},
  guard/.style={draw=tpggray, rounded corners=1.5pt, fill=black!3, align=center, inner sep=2.2mm}
]
  \node[phase, draw=tpgblue, fill=tpgblue!9, text width=32mm] (train) {
    \textbf{Train}\\2020--2023
  };
  \node[phase, draw=tpgteal, fill=tpgteal!9, text width=20mm, right=1.5mm of train] (val) {
    \textbf{Validate}\\2024
  };
  \node[phase, draw=tpgorange, fill=tpgorange!9, text width=27mm, right=1.5mm of val] (test) {
    \textbf{Final test}\\2025--2026
  };

  \draw[-{Stealth[length=2mm]}, very thick, draw=tpggray]
    ([yshift=8mm]train.west) -- ([yshift=8mm]test.east)
    node[midway, above=0.7mm, font=\sffamily\scriptsize] {chronological allocation};

  \node[guard, text width=84mm, below=4mm of val] (lock) {
    \textbf{Group lock before allocation:}
    event identity $\mid$ near-duplicate claim $\mid$ perceptually identical image
  };
  \draw[-{Stealth[length=1.8mm]}, semithick, draw=tpggray] (lock.north -| train) -- (train.south);
  \draw[-{Stealth[length=1.8mm]}, semithick, draw=tpggray] (lock.north -| val) -- (val.south);
  \draw[-{Stealth[length=1.8mm]}, semithick, draw=tpggray] (lock.north -| test) -- (test.south);

\end{tikzpicture}
}
\caption{Primary temporal split and leakage controls. A locked event, near-duplicate claim, or reused image cannot cross partitions.}
\label{fig:temporal-protocol}
\end{figure}

This study treats \datasetname{} as a benchmark only if the filtered corpus contains at least 2,000 cases, 300 cases per label, 50 test events, and two independent evidence clusters for at least 70\% of retained claims. If the collection misses any threshold, it is used as a temporal and provenance stress test instead.

